% !TeX root = ../教材题库_全部条目_转写工作稿.tex
% Source: XJ-B1-C04-RV-P152-Y0261
% 第4章 / 章末复习 复习题四 / 教材 p.152 / PDF 第159页
% OCR未校对：公式与文字均可能有误；未完成数学TeX转写；禁止直接发给学生
% 原题及相邻原解截图（核对用）：../题目原图/XJ-B1-C04-RV-P152-Y0261-01.webp
% 本文件只在转写完成后作为题干使用；OCR 条目尚不是可靠题干。
\begin{bankitem}{复习题四 第11题}{第4章 / 章末复习 复习题四 / 教材 p.152 / PDF 第159页}{XJ-B1-C04-RV-P152-Y0261}{OCR未校对：公式与文字均可能有误；未完成数学TeX转写；禁止直接发给学生}
\begin{ocrdraft}
[教材p.152 / PDF第159页]
\par
11.根据有关资料，围棋状态空间复杂度的上限M约为33丛而可观测宇宙
\par
中普通物质的原子总数N约为1080，已知lg30．4\&从数量级的角度考虑，下
\par
列各数中与最接近的是（）
\par
（A)1033
\par
(B)1053
\par
(C)1073
\par
（D)109
\end{ocrdraft}
\end{bankitem}
